\documentclass[conference]{IEEEtran} \IEEEoverridecommandlockouts


\usepackage{amsmath,amssymb,amsfonts} \usepackage{algorithmic} \usepackage{algorithm} \usepackage{graphicx} \usepackage{textcomp} \usepackage{xcolor} \usepackage{booktabs} \usepackage{multirow} \usepackage{url}


\def\BibTeX{{\rm B\kern-.05em{\sc i\kern-.025em b}\kern-.08em     T\kern-.1667em\lower.7ex\hbox{E}\kern-.125emX}}

\begin{document}




\title{LP-Pruned Quantum Branch-and-Bound for the 0/1 Knapsack Problem: Query-Complexity Analysis and Spectral Validation}

\author{ \IEEEauthorblockN{ Shiva Sai$^{1}$, Kumudini$^{1}$, Purusharth$^{1}$, Rakshit$^{2}$, Veerabhadra$^{2}$, C. Thirmal$^{2}$ }

\IEEEauthorblockA{$^{1}$Student, IIT Bachupally, India}

\IEEEauthorblockA{$^{2}$Professor, Harvard University, USA} }

\maketitle




\begin{abstract} The 0/1 knapsack problem is a fundamental NP-hard combinatorial optimization problem with applications spanning resource allocation, cryptography, and logistics. The best provable quantum algorithms either match classical meet-in-the-middle at $O(2^{n/2})$ or rely on unproven heuristic assumptions (Heuristic~2) and infeasible QRAQM hardware. We present a hybrid algorithm for knapsack that couples LP-guided reduced-cost variable fixing with Montanaro's quantum branch-and-bound framework. Unlike prior quantum knapsack methods---which are generic, heuristic, or QRAQM-dependent---our method is simultaneously provable, gate-model feasible (QRAQM-free), knapsack-specific, and instance-adaptive. Our three-phase hybrid algorithm uses: (1)~reduced-cost variable fixing to provably eliminate items from the search space, reducing the problem to $m = \alpha n$ core items where $\alpha \in [0, 1]$ depends on instance structure; (2)~a classical threshold-search outer loop invoking quantum detection, with composed query complexity $O\!\left(\sqrt{T_{\max}}\,m\log m\log V_{\max}\right)$, where $T_{\max}$ is the maximum LP-pruned tree size encountered over the tested thresholds; and (3)~classical post-processing to reconstruct the global optimum. We prove instance-dependent complexity bounds showing that for weakly constrained instances, LP preprocessing reduces the effective problem size to near-constant, while for strongly constrained instances with substantially larger LP-pruned trees, the quantum phase achieves a quadratic query reduction relative to classical traversal of the same backtracking tree under the Montanaro oracle-model assumptions. Beyond complexity counting, we computationally test the spectral conditions for the quantum phase on the LP-pruned tree. We verify the effective spectral-gap inequality and phase-zero weight bound, finding perfect agreement with the analytical detection criterion on 889 completed simulated trees ($T_{\text{LP}} \le 2000$). We benchmark against strong classical solvers (dynamic programming, Google OR-Tools), clarifying that the query-complexity reduction is relevant only in the large-coefficient, LP-hard regime where DP is infeasible and branch-and-bound explodes. A fault-tolerant resource estimate for our answer-agnostic reversible oracle ($31$--$45$ logical qubits) establishes baseline component costs for the predicate $V$. Finally, we characterize the NISQ noise boundary of amplitude amplification on IBM's \texttt{ibm\_fez} processor (156-qubit Heron~r2) for core sizes $n \in \{3, 4, 5, 6\}$. We do not claim wall-clock quantum speedup, hardware advantage, or complete end-to-end full-walk gate complexity reduction. \end{abstract}

\begin{IEEEkeywords} quantum computing, knapsack problem, branch-and-bound, quantum walk, LP relaxation, variable fixing, combinatorial optimization \end{IEEEkeywords}




\section{Introduction}

The 0/1 knapsack problem is one of the most studied problems in combinatorial optimization: given $n$ items, each with a weight $w_i$ and a value $v_i$, and a capacity $W$, find the subset that maximizes total value without exceeding the capacity. Despite its simple formulation, the problem is NP-hard, and the best classical exact algorithm---the Horowitz--Sahni meet-in-the-middle approach---runs in $O(2^{n/2})$ time and space \cite{ref2}.

Quantum computing has long promised speedups for combinatorial problems. Grover's algorithm \cite{ref8} provides a quadratic speedup for unstructured search, and its generalization via amplitude amplification \cite{ref11} is a powerful tool. However, applying Grover search directly to knapsack yields $O(2^{n/2})$---no improvement over classical meet-in-the-middle.

More sophisticated quantum approaches based on quantum walks have achieved better asymptotic complexities. Bernstein et al. \cite{ref16} achieved $O(2^{0.241n})$, and Helm and May \cite{ref17} further improved this to $O(2^{0.226n})$. However, these results rely on Heuristic~2---an unproven assumption about quantum walk update times---which Bonnetain et al. \cite{bonnetain2020improved} have shown to be mathematically unsound in general. Furthermore, these algorithms require QRAQM (Quantum Random Access Memory with Quantum Addressing), a hardware model that does not exist and faces severe engineering challenges \cite{ref18}.

This creates a gap in the literature: can LP structure be combined with provable quantum tree-search query bounds without unproven heuristics or QRAQM?

\subsection{Contributions and Novelty}

Our algorithm uniquely combines exact LP reduced-cost fixing with a quantum tree search. Our defensible contributions are:

\begin{enumerate}
    \item Exact LP/reduced-cost variable fixing for reducing the knapsack core before quantum tree search.
    \item A threshold-optimization composition using a classical binary-search outer loop and a Montanaro-style quantum detection primitive.
    \item A composed query bound $O\!\left(\sqrt{T_{\max}}\,m\log m\log V_{\max}\right)$, with $T_{\max}$ explicitly defined as the maximum threshold-dependent LP-pruned tree size encountered by the outer loop.
    \item Computational spectral validation of the implemented walk on 889 completed simulated trees with $T_{\text{LP}} \le 2000$, including effective spectral-gap inequality, phase-zero marked weight, and analytical phase-estimation acceptance separation.
    \item Small-scale IBM hardware characterization of the reversible threshold predicate/core circuit, explicitly limited to the executed small instances.
\end{enumerate}

We do not claim wall-clock quantum acceleration or a complete gate-complexity bound for the full LP-bounded walk, because the reversible LP heuristic $h$ is not synthesized.

The remainder of this paper is organized as follows. Section~II reviews related work. Section~III establishes mathematical preliminaries. Section~IV presents our algorithm. Section~V provides the theoretical analysis. Section~VI describes our experimental methodology and results. Section~VII discusses implications and limitations. Section~VIII concludes.

% ==============================================================
% II. LITERATURE REVIEW
% ==============================================================
\section{Literature Review}

\textit{Classical exact methods.} Dantzig's LP relaxation yields a closed-form fractional optimum and the tight Dantzig upper bound that remains the workhorse of exact knapsack solvers \cite{ref1}. The Horowitz--Sahni meet-in-the-middle (MitM) algorithm splits the items into two halves and merges sorted partial sums, attaining $O(2^{n/2})$ time and space \cite{ref2}---still the best \emph{provable} worst-case bound for exact 0/1 knapsack and the natural classical target for any quantum method. Kellerer, Pferschy, and Pisinger consolidate dynamic programming (pseudo-polynomial $O(nW)$), greedy bounds, and branch-and-bound (B\&B) \cite{ref3}. Pisinger's \emph{core} concept is central to our approach: empirically only items whose efficiency lies near the LP split affect the optimum \cite{ref4}, and reduced-cost variable fixing turns this observation into an \emph{exact} reduction. Closely related subset-sum work includes the representation technique of Becker, Coron, and Joux \cite{ref5} and its refinements \cite{ref6}, while Plateau and Elkihel analyze packing-tree pruning that prefigures modern B\&B \cite{ref7}.

\textit{Quantum search and its limits.} Grover's algorithm searches an unstructured space of size $N$ in $O(\sqrt{N})$ \cite{ref8}, later sharpened to exact zero-failure search \cite{ref10} and generalized to amplitude amplification \cite{ref11}, with optimization variants via pure adaptive search \cite{ref12} and D\"urr--H{\o}yer minimum finding \cite{ref13}. Applied naively to knapsack, however, Grover yields only $O(2^{n/2})$---no asymptotic gain over MitM---and Bennett et al.\ proved that black-box quadratic speedups cannot break the $2^{n/2}$ barrier for NP-complete problems absent additional structure \cite{ref9}. Ambainis' adversary method supplies matching query lower bounds \cite{ref34}, framing precisely what a structure-exploiting method must overcome.

\textit{Quantum walks and the heuristic/QRAQM barrier.} Structured search improvements arise from quantum walks \cite{ref14}, formalized by the MNRS framework \cite{ref15}. For subset-sum, Bernstein et al.\ reach $O(2^{0.241n})$ \cite{ref16} and Helm--May $O(2^{0.226n})$ \cite{ref17}---the strongest asymptotics known. Yet two caveats motivate our work: (i)~they assume \emph{Heuristic~2}, an unproven claim about walk-update costs that Bonnetain et al.\ showed to be mathematically unsound in general, recovering provable but weaker bounds \cite{bonnetain2020improved}; and (ii)~they require QRAQM (quantum-addressable quantum memory), whose feasibility is heavily contested \cite{ref18,ref19}. The best ``quantum'' knapsack asymptotics are therefore simultaneously unproven and hardware-infeasible.

\textit{Implementable quantum knapsack and hybrids.} A second line prioritizes constructibility on gate-model hardware. Grover-based knapsack solvers \cite{ref20,ref21} and quadratic/multidimensional extensions \cite{ref22} are implementable but do not surpass MitM. Variational methods---QAOA \cite{ref25}, Grover adaptive search \cite{ref26}, and energy-systems applications \cite{ref27}---run on NISQ devices but lack performance guarantees. Quantum-inspired metaheuristics (genetic \cite{ref28}, wolf-pack \cite{ref29}, harmonic-oscillator \cite{ref30}) report practical gains without provable query reduction. The most relevant provable, QRAQM-free tool is Montanaro's query reduction of B\&B \cite{ref24}, building on quantum backtracking \cite{ref35}; Sanavio et al.\ apply hybrid B\&B to integer linear programs \cite{ref36}, and Wilkening et al.'s Quantum Tree Generator (QTG) builds superpositions of feasible knapsack solutions for QAOA \cite{ref23}. Hardware context is set by Preskill's NISQ analysis \cite{ref31}, shallow-circuit capabilities \cite{ref32}, and the Qiskit toolchain \cite{ref33}.

\textit{Closest prior work and our position.} Three works are nearest to ours. Montanaro's quantum B\&B \cite{ref24} is provable and QRAQM-free but \emph{generic}: it bounds a quantum walk on an abstract tree of size $T$ by $\tilde{O}(\sqrt{T}\,d)$ without analyzing any specific problem, the role of LP bounds, or the induced tree size for knapsack. The QTG of Wilkening et al.\ \cite{ref23} is knapsack-specific and gate-model, but feeds feasible-solution superpositions to a \emph{variational} (QAOA) routine and performs no LP-based variable fixing, inheriting QAOA's lack of guarantees. Bonnetain et al.\ \cite{bonnetain2020improved} restore provability to walk-based subset-sum yet still assume QRAQM. Our algorithm uniquely combines: (a)~\emph{exact} LP reduced-cost fixing with Montanaro's provable, QRAQM-free quantum B\&B specifically for knapsack, (b)~an analysis of the resulting LP-pruned tree size $T_{\text{LP}} \le T$, and (c)~computational spectral validation of the explicit walk operator on completed simulable trees. Table~\ref{tab:related} contrasts these properties.

\begin{table}[htbp]
\caption{Capability comparison with the closest prior quantum methods. The lower block lists capabilities that, to our knowledge, no prior method provides simultaneously.}
\label{tab:related}
\begin{center}
\footnotesize
\setlength{\tabcolsep}{4pt}
\begin{tabular}{lcccc}
\toprule
\textbf{Capability} & \textbf{Mont.} & \textbf{QTG} & \textbf{Bonn.} & \textbf{Ours} \\
 & \cite{ref24} & \cite{ref23} & \cite{bonnetain2020improved} & \\
\midrule
Provable query reduction               & Yes & No\textsuperscript{\S} & Yes & \textbf{Yes} \\
Gate-model feasible (no QRAQM)         & Yes & Yes & No  & \textbf{Yes} \\
Knapsack-specific                      & No  & Yes & No  & \textbf{Yes} \\
\midrule
Exact LP variable fixing               & No  & No  & No  & \textbf{Yes} \\
Instance-adaptive$^\dagger$            & No  & No  & No  & \textbf{Yes} \\
Explicit walk spectral validation$^\ddagger$       & No  & No  & No  & \textbf{Yes} \\
Hardware demonstration                 & No  & No  & No  & \textbf{Yes} \\
\bottomrule
\multicolumn{5}{l}{\footnotesize Mont.: Montanaro B\&B; QTG: Quantum Tree Generator; Bonn.: Bonnetain et al.} \\
\multicolumn{5}{l}{\footnotesize $^\dagger$Query bound depends on the LP-pruned tree sizes encountered.}\\
\multicolumn{5}{l}{\footnotesize $^\ddagger$Explicit walk spectra and analytical detection criteria are evaluated computationally on completed simulable trees.} \\
\multicolumn{5}{l}{\footnotesize \textsuperscript{\S}Variational (QAOA): no performance guarantee.} \\
\end{tabular}
\end{center}
\end{table}

The upper block of Table~\ref{tab:related} shows that no single prior method is simultaneously provable, gate-model feasible, and knapsack-specific. The lower block isolates the capability our work contributes: an algorithm that couples \emph{exact} LP variable fixing to a provable, QRAQM-free quantum walk, is \emph{instance-adaptive}---with query complexity adapting to the LP-pruned tree sizes encountered---and whose spectral conditions and analytical detection criterion are computationally tested rather than assumed. This combination is what distinguishes the present work, and it is what we substantiate empirically in Section~\ref{sec:experiments}.

% ==============================================================
% III. PRELIMINARIES
% ==============================================================
\section{Preliminaries}

\subsection{Problem Definition}

The 0/1 Knapsack Problem is defined as follows. Given a set of $n$ items, where item $i$ has weight $w_i \in \mathbb{Z}^+$ and value $v_i \in \mathbb{Z}^+$, and a capacity $W \in \mathbb{Z}^+$, find a binary vector $x^* = (x_1^*, x_2^*, \ldots, x_n^*) \in \{0,1\}^n$ that solves
\begin{equation}
    \max \sum_{i=1}^{n} v_i x_i \quad \text{subject to} \quad \sum_{i=1}^{n} w_i x_i \leq W.
    \label{eq:knapsack}
\end{equation}
We denote the optimal value by $Z^* = \sum_{i} v_i x_i^*$. Without loss of generality, we assume $w_i \leq W$ for all $i$ (items heavier than the capacity are trivially excluded) and that all efficiencies $e_i = v_i / w_i$ are distinct.

\subsection{LP Relaxation, Split Item, and Dantzig Bound}

The LP relaxation of \eqref{eq:knapsack} permits fractional assignments $x_i \in [0,1]$. Its solution has a closed-form structure due to Dantzig \cite{ref1}.

\textbf{Definition 1 (Efficiency Ordering).} Order items so that $e_1 > e_2 > \cdots > e_n$ where $e_i = v_i / w_i$.

\textbf{Definition 2 (Split Item).} The \textit{split item} $s$ is the smallest index such that
\begin{equation}
    \sum_{i=1}^{s} w_i > W \quad \text{and} \quad \sum_{i=1}^{s-1} w_i \leq W.
    \label{eq:split}
\end{equation}

The LP optimal solution is then
\begin{equation}
    x_i^{LP} = \begin{cases}
    1 & \text{if } i < s, \\
    (W - \sum_{j=1}^{s-1} w_j) / w_s & \text{if } i = s, \\
    0 & \text{if } i > s,
    \end{cases}
    \label{eq:lpsol}
\end{equation}
and the \textit{Dantzig bound} is $Z_{LP} = \sum_{i=1}^{n} v_i x_i^{LP}$. This is computable in $O(n)$ time after sorting, and satisfies $Z^* \leq Z_{LP}$.

\textbf{Definition 3 (Core Items).} Given a parameter $\delta > 0$, the \textit{core} $C$ is the set of items with efficiencies close to the split item:
\begin{equation}
    C = \{i : |e_i - e_s| \leq \delta \cdot e_s\}.
    \label{eq:core}
\end{equation}
The remaining items are partitioned into \textit{fixed-in} items $F_1 = \{i : e_i > e_s + \delta \cdot e_s\}$ and \textit{fixed-out} items $F_0 = \{i : e_i < e_s - \delta \cdot e_s\}$.

The core concept was introduced by Pisinger \cite{ref4}, who showed empirically that $|C| \ll n$ for typical instances. Our algorithm exploits this structure theoretically.

\subsection{Amplitude Amplification}

We use the following standard result \cite{ref11}.

\textbf{Theorem (Brassard et al., 2002).} Let $\mathcal{A}$ be a quantum algorithm that produces $\mathcal{A}|0\rangle = \sin\theta\,|\Psi_{\text{good}}\rangle + \cos\theta\,|\Psi_{\text{bad}}\rangle$. Then amplitude amplification finds a good state in $O(1/\sqrt{p})$ oracle calls where $p = \sin^2\theta$.

\subsection{Quantum Branch-and-Bound (Montanaro)}

Montanaro~\cite{ref24} proved that quantum computers can speed up classical branch-and-bound algorithms:

\textbf{Theorem (Montanaro, 2020).} Let $\mathcal{B}$ be a classical B\&B algorithm that determines if there is a valid assignment evaluating to at least $\tau$, exploring a tree of size $T$ and maximum depth $d$. There is a quantum algorithm that solves the same detection problem with failure probability $\le 1/3$ using
\begin{equation}
    O(\sqrt{T}\, d \log d)
    \label{eq:montanaro}
\end{equation}
queries to the branching and bounding heuristics.

The algorithm uses a quantum walk on the implicit B\&B tree and does not require QRAQM---only the ability to evaluate the bounding function (our LP relaxation) as a quantum oracle. This is the key quantum ingredient in our approach.

\subsection{Reduced-Cost Variable Fixing}

Given a lower bound $Z_{\text{LB}}$ (e.g., from a greedy solution), an item~$i$ with $x_i^{LP} = 1$ can be \textit{provably fixed} to~1 if removing it causes the LP upper bound to drop below $Z_{\text{LB}}$. Similarly, an item~$j$ with $x_j^{LP} = 0$ can be fixed to~0 if forcing it in causes the LP bound to drop below $Z_{\text{LB}}$. This technique, due to Pisinger~\cite{ref4}, is exact---no heuristic assumptions are involved.

\subsection{Meet-in-the-Middle Framework}

The Horowitz--Sahni algorithm \cite{ref2} partitions the $n$ items into two disjoint sets $A$ and $B$ with $|A| = |B| = n/2$. It enumerates all $2^{n/2}$ subsets of $A$, storing pairs $(w(S_A), v(S_A))$ for each subset $S_A \subseteq A$, where $w(S_A) = \sum_{i \in S_A} w_i$ and $v(S_A) = \sum_{i \in S_A} v_i$. It does the same for $B$. Then, for each subset $S_A$ of $A$, it finds the best compatible subset $S_B$ of $B$ satisfying $w(S_A) + w(S_B) \leq W$ using binary search on the sorted list of $B$-subsets. The total time and space are both $O(2^{n/2})$.

% ==============================================================
% IV. PROPOSED ALGORITHM
% ==============================================================
\section{Proposed Algorithm}

Our hybrid algorithm operates in three phases: classical LP preprocessing, quantum branch-and-bound over the reduced core, and classical post-processing.

\subsection{Phase 1: LP Preprocessing via Reduced-Cost Variable Fixing}

The classical preprocessing exploits the LP relaxation and reduced-cost analysis to provably partition items into three disjoint sets.

\begin{algorithm}[h]
\caption{Reduced-Cost Variable Fixing}
\label{alg:filter}
\begin{algorithmic}[1]
\REQUIRE Items $(w_i, v_i)_{i=1}^n$, capacity $W$
\ENSURE Partition $(F_1, C, F_0)$, residual capacity $W'$
\STATE Compute efficiencies $e_i \leftarrow v_i / w_i$; sort decreasingly
\STATE Compute greedy lower bound $Z_{\text{LB}}$
\STATE Compute Dantzig upper bound $Z_{\text{LP}}$
\FOR{each item $i \in \{1, \ldots, n\}$}
    \IF{$x_i^{LP} = 1$ (item before split)}
        \STATE $Z'_i \leftarrow$ LP bound with item $i$ forced out
        \IF{$Z'_i < Z_{\text{LB}}$}
            \STATE $F_1 \leftarrow F_1 \cup \{i\}$ \COMMENT{Provably in optimal}
        \ELSE
            \STATE $C \leftarrow C \cup \{i\}$
        \ENDIF
    \ELSE
        \STATE $Z'_i \leftarrow$ LP bound with item $i$ forced in
        \IF{$Z'_i < Z_{\text{LB}}$}
            \STATE $F_0 \leftarrow F_0 \cup \{i\}$ \COMMENT{Provably out of optimal}
        \ELSE
            \STATE $C \leftarrow C \cup \{i\}$
        \ENDIF
    \ENDIF
\ENDFOR
\STATE $W' \leftarrow W - \sum_{i \in F_1} w_i$; $V_1 \leftarrow \sum_{i \in F_1} v_i$
\RETURN $(F_1, C, F_0)$, $W'$, $V_1$
\end{algorithmic}
\end{algorithm}

Unlike heuristic delta-based filtering, reduced-cost fixing is \textit{provably exact}: no item in $F_1$ or $F_0$ can change the optimal solution. Let $m = |C|$ and $\alpha = m/n$.

\subsection{Phase 2: Quantum Branch-and-Bound on Core}

The core contribution of this paper. We construct a classical B\&B tree over the $m$ core items using LP relaxation bounds at each node, then apply Montanaro's quantum walk~\cite{ref24} to search this tree in $$O\!\left(\sqrt{T_{\max}}\,m\log m\log V_{\max}\right)$$ queries.

\subsubsection{Classical B\&B Tree Structure}

At each node of the B\&B tree, we branch on a core item $c_i$ (include/exclude) and compute the LP relaxation bound on the remaining subproblem. A node is pruned when its LP bound falls below the current best known integer solution.

\textbf{Bounding oracle.} The LP relaxation for the residual knapsack at each B\&B node is computable in $O(m)$ time classically. The LP relaxation defines the bounding heuristic $h$ required by the abstract query model.

\subsubsection{Reversible LP-Bound Oracle ($h$)}

To implement the bounding heuristic $h$ as a quantum oracle, we construct a fixed-precision integer circuit that evaluates the exact threshold rule from Lemma \ref{lem:precision}. The fixed-precision calculation evaluates a scaled integer bound $\tilde{Z}_{scaled} = V_{int} 2^k + W_{res} \lfloor v_s 2^k / w_s \rfloor$. Because $\tilde{Z}_{scaled}$ is guaranteed to be an integer, the condition $\tilde{Z}'_i \le Z_{LB} - 1/W_{max}$ is evaluated by comparing $\tilde{Z}_{scaled}$ against a classically pre-computed threshold $T = \lfloor Z_{LB} 2^k - 2^k/W_{max} \rfloor$ without any precision loss.

The register layout allocates quantum memory for the scaled dividend, divisor, quotient, remainder, residual capacity, and the intermediate additions and multiplications. The arithmetic relies on the Thapliyal restoring divider \cite{thapliyal2016} for calculating the fixed-point efficiency, the Vedral-Barenco-Ekert (VBE) multiplier \cite{vedral1996} for computing the residual product, and the zero-ancilla Draper QFT adder \cite{draper2000}. By mapping the division remainder to the ancillary register required by Thapliyal's in-place division, the ancillary overhead is strictly bounded. Resource estimates for this layout are provided in Table~\ref{tab:h_resources}. Note that while the exact logical and ancillary qubit counts are established, full gate-level decomposition of the arithmetic blocks is deferred.

\begin{table}[ht]
\centering
\caption{Resource estimates for the LP-bound bounding oracle ($h$) at varying core sizes ($m$). The $(m, W_{max})$ pairings shown (15/31/63) are representative illustrative example values chosen for the demonstration, not a derived or assumed functional relationship between core size and maximum item weight --- these are independent parameters in general knapsack instances. Qubit counts are derived from the register layout mapping to the Draper, Vedral-Barenco-Ekert, and Thapliyal reversible arithmetic blocks. The fractional fixed-point precision ($k$) scales with $W_{max}$ per Lemma \ref{lem:precision}. Depth, gate, and T-counts are not yet available as full gate-level decomposition has been deferred for this predicate.}


We have subsequently synthesized and tested the gate-level circuit specifically for the $m=4$ ($W_{max}=15$) case, providing a non-deferred evaluation of oracle $h$. Against a stratified sample of 25 classical benchmark instances (representing varying hardness classes), the custom-simulated gate-level circuit achieved 0 mismatches against standard integer arithmetic, and passed an injected-fault unit test. The measured hardware metrics for $m=4$ yield 96 qubits, 8,170 gates, and a depth of 3,844. For the T-count, a naive estimate weighting Toffoli and controlled-controlled-phase (`ccp`) gates uniformly at 7 T-gates yields 17,633. However, utilizing the standard decomposition by Barenco et al.~\cite{barenco1995} (Lemma 6.1) for arbitrary doubly-controlled $C^2P(\theta)$ unitaries, which uses three controlled-phase rotations, and synthesizing these continuous rotations via the Ross-Selinger algorithm at precision $\epsilon=10^{-10}$, the T-count scales to roughly $\approx 1.2 \times 10^6$. We report both the naive and Ross-Selinger-scaled figures to distinguish the raw gate counts from the fault-tolerant synthesis penalty. Note that $m=6$ and $m=8$ remain at the analytical estimate status, as full synthesis was not extended to those larger core sizes.

\label{tab:h_resources}
\begin{tabular}{ccccccc}
\hline\hline
$m$ ($W_{max}$) & Precision $k$ & Logical Qubits & Ancilla Qubits & Total Qubits & Depth / Gates & T-count \\
\hline
4 (15) (Analytical) & 8 & \multicolumn{3}{c}{\textit{(Analytical estimate: 82 qubits)}} & \multicolumn{2}{c}{\textit{N/A --- gate-level decomposition deferred}} \\
4 (15) (MEASURED) & 8 & 84 & 12 & 96 & 3,844 / 8,170 & $\approx 1.2\times 10^6$$^\ddagger$ (17,633$^\star$) \\
6 (31) & 10 & 101 & 1 & 102 & N/A --- gate-level decomposition deferred & N/A \\
8 (63) & 12 & 121 & 1 & 122 & N/A --- gate-level decomposition deferred & N/A \\
\hline\hline
\multicolumn{7}{l}{\footnotesize $^\ddagger$Ross-Selinger estimate $\epsilon=10^{-10}$ using Barenco et al.~\cite{barenco1995} 3-rotation decomposition.} \\
\multicolumn{7}{l}{\footnotesize $^\star$Naive estimate treating `ccp` gates as standard Toffolis (7 T-gates).} \\
\end{tabular}
\end{table}

\subsubsection{Quantum Walk Application}

We apply Montanaro's quantum B\&B framework~\cite{ref24} to the implicit tree defined by our LP-bounded branching strategy:

\begin{algorithm}[h]
\caption{Classical Threshold-Search Outer Loop (Phase 2)}
\label{alg:qbb}
\begin{algorithmic}[1]
\REQUIRE Core items $C$ with $|C| = m$, residual capacity $W'$, lower bound $V_1$
\ENSURE Optimal core assignment $x^*_C$
\STATE Initialize binary search interval for threshold $\tau$
\WHILE{search interval $> 0$}
\STATE Define LP-pruned fixed-threshold tree for $\tau$
\STATE Invoke quantum detection on that tree
\STATE Update binary-search interval from the detection result
\ENDWHILE
\STATE Recover the optimum threshold/value after the loop
\RETURN Optimal core assignment $x^*_C$
\end{algorithmic}
\end{algorithm}

For each tested threshold $\tau_i$, the fixed-threshold detection primitive requires $O(\sqrt{T_{\tau_i}}\,m\log m)$ queries under the Montanaro oracle-model assumptions. Since the classical outer loop tests at most $O(\log V_{\max})$ thresholds, summing these calls gives $O(\sqrt{T_{\max}}\,m\log m\log V_{\max})$ queries, where $T_{\max}=\max_i T_{\tau_i}$.

\subsubsection{Reversible Knapsack Oracle (Answer-Agnostic)}
\label{sec:oracle-impl}

The quantum subroutines require a constructible oracle that marks states satisfying the threshold predicate $\text{feasible}(x) \wedge \text{value}(x) \geq \tau$, where $\tau$ is the binary-search threshold---\textit{never} the (unknown) optimal value. We implement this as a genuine gate-level reversible circuit using in-place Draper (QFT) arithmetic with \textit{zero ancilla} beyond two accumulator registers: (1)~accumulate $\sum_i w_i x_i$ into a weight register and subtract $W{+}1$, so its two's-complement sign bit equals feasibility; (2)~accumulate $\sum_i v_i x_i$ and subtract $\tau$, so its sign bit encodes $\text{value} \geq \tau$; (3)~flip the phase conditioned on both flags; (4)~uncompute. 

Table~\ref{tab:oracle} details the exact logical resource requirements for this evaluation/threshold predicate $V$. Crucially, we synthesize the predicate $V$ only; the LP bounding heuristic $h$ is not synthesized to a reversible Clifford+T circuit. 

\begin{table}[htbp]
\caption{Exact Logical Predicate Resources and Analytical T-Count Estimates}
\label{tab:oracle}
\begin{center}
\footnotesize
\begin{tabular}{lccccc}
\toprule
\textbf{Core Size ($m$)} & \textbf{Qubits} & \textbf{Depth} & \textbf{Total Gates} & \textbf{2Q Gates} & \textbf{T-Count Estimate$^\dagger$} \\
\midrule
4 & 19 & 266 & 1353 & 573 & $\approx 188{,}004$ \\
6 & 22 & 300 & 1788 & 773 & $\approx 227{,}631$ \\
8 & 25 & 327 & 2243 & 977 & $\approx 250{,}404$ \\
\bottomrule
\multicolumn{6}{l}{\footnotesize $^\dagger$Heuristic analytical estimate assuming Ross-Selinger arbitrary $R_z$ synthesis at $\epsilon = 10^{-10}$.} \\
\multicolumn{6}{l}{\footnotesize The LP heuristic $h$ was NOT synthesized.}
\end{tabular}
\end{center}
\end{table}

This construction uses only the problem data and $\tau$, so the optimum is \textit{discovered} by the search rather than supplied to it---unlike a diagonal oracle built from a precomputed solution. We verified by exhaustive statevector simulation that, for all thresholds, the oracle marks exactly the brute-force-correct set and returns ancillae to $|0\rangle$, and that one amplitude-amplification stage with the full reversible circuit matches the verified phase action to $<10^{-13}$. In a separate small-scale component experiment, a D\"urr--H{\o}yer-style maximum-finding loop built on this oracle recovered the exact optimum on the tested $n \le 5$ statevector instances. This experiment is not the optimization loop used by quantum\_bb\_optimize.

\subsection{Phase 3: Classical Post-Processing}

\begin{algorithm}[h]
\caption{Classical Reconstruction}
\label{alg:postprocess}
\begin{algorithmic}[1]
\REQUIRE $(F_1, C, F_0)$, quantum solution $x^*_C$, capacity $W$
\ENSURE Globally optimal solution $x^*$
\STATE $x^* \leftarrow F_1 \cup \{i \in C : x^*_{C,i} = 1\}$
\STATE $Z^* \leftarrow \sum_{i \in x^*} v_i$
\RETURN $x^*$, $Z^*$
\end{algorithmic}
\end{algorithm}

Since reduced-cost fixing is exact, the combined solution $F_1 \cup x^*_C$ is globally optimal. No swap verification is needed.

\subsection{Overall Algorithm}

\begin{algorithm}[h]
\caption{LP-Pruned Quantum B\&B Knapsack Solver}
\label{alg:hybrid}
\begin{algorithmic}[1]
\REQUIRE Items $(w_i, v_i)_{i=1}^n$, capacity $W$
\ENSURE Optimal solution $x^*$ and value $Z^*$
\STATE \textbf{Phase 1:} $(F_1, C, F_0), W', V_1 \leftarrow$ \textsc{RC-Fix}$(w, v, W)$ \COMMENT{Alg.~\ref{alg:filter}}
\STATE $m \leftarrow |C|$
\IF{$m = 0$}
    \RETURN $x^* = F_1$, $Z^* = V_1$
\ELSIF{$m \leq m_{\text{thresh}}$}
    \STATE Solve core classically via B\&B
\ELSE
    \STATE \textbf{Phase 2:} $x^*_C \leftarrow$ \textsc{Quantum-B\&B}$(C, W', V_1)$ \COMMENT{Alg.~\ref{alg:qbb}}
\ENDIF
\STATE \textbf{Phase 3:} $x^* \leftarrow F_1 \cup x^*_C$
\RETURN $x^*$, $Z^* = V_1 + v_C(x^*_C)$
\end{algorithmic}
\end{algorithm}

\subsection{Algorithmic Extensions}



\subsubsection{Multidimensional Knapsack Extension}
Our framework extends cleanly to the $d$-dimensional 0/1 knapsack problem. We generalize the LP fixing by using a surrogate relaxation: the minimum of $d$ single-constraint Dantzig bounds. Sorting items by per-constraint efficiencies $v_i/w_{ij}$ yields provably valid upper bounds for reduced-cost fixing (because each single-constraint Dantzig bound maximizes the objective over a strict superset of the full $d$-dimensional feasible region, making their minimum a valid relaxation). Table~\ref{tab:multidim} summarizes the empirical behavior of $\alpha$ across 20 independent trials per configuration for strongly correlated instances ($n=24$). For these multidimensional instances, weights are drawn independently as $w_{ij} \sim U(1, 100)$, and values are tied to the average weight via $v_i = \lfloor \frac{1}{d} \sum_{j=1}^d w_{ij} \rfloor + 10$. The results show that strongly correlated multidimensional instances heavily resist classical variable fixing, saturating to $\alpha = 1.0$ instantly at $d \ge 2$, meaning zero variables are fixed by Phase 1. The consequences for quantum query complexity require a separate analysis of the resulting pruned trees.

\begin{table}[htbp]
\caption{Core size ratio $\alpha$ vs. dimensionality $d$ for strongly correlated instances ($n=24$, 20 trials)}
\label{tab:multidim}
\begin{center}
\begin{tabular}{lcccc}
\toprule
\textbf{Dimensionality} & $d=1$ & $d=2$ & $d=3$ & $d=4$ \\
\midrule
\textbf{Core Ratio $\alpha$ (Mean $\pm$ 95\% CI)} & $0.96 \pm 0.06$ & $1.00 \pm 0.00$ & $1.00 \pm 0.00$ & $1.00 \pm 0.00$ \\
\bottomrule
\end{tabular}
\end{center}
\end{table}

% ==============================================================
% V. THEORETICAL ANALYSIS
% ==============================================================
\section{Theoretical Analysis}

\subsection{Correctness}

\textbf{Lemma 1a (Precision requirement for robust reduced-cost fixing).} \label{lem:precision} \textit{Assume a 0/1 knapsack instance with capacity $W$, max weight $W_{max}$, and a greedy integer lower bound $Z_{LB}$. Let $Z'_i$ be the exact perturbed LP bound for item $i$.
If the LP bound is computed using $k$-bit floor truncation for the split item's efficiency $e_s$, yielding a truncated bound $\tilde{Z}'_i$, the quantum detection loop will produce the \textbf{exact same} variable fixing partition $(F_1, C, F_0)$ as infinite-precision classical Algorithm 1, provided two conditions are met:
\begin{enumerate}
    \item The classical check $Z'_i < Z_{LB}$ is replaced by the robust check $\tilde{Z}'_i \le Z_{LB} - 1/W_{max}$.
    \item The precision satisfies $k > 2 \log_2(W_{max})$.
\end{enumerate}}

\textit{Note:} The adoption of the robust check $\tilde{Z}'_i \le Z_{LB} - 1/W_{max}$ entirely eliminates the exception for degenerate instances ($Z'_i = Z_{LB}$). Because the maximum truncation error $W_{max} 2^{-k}$ is strictly less than the safety margin $1/W_{max}$, the truncated bound will never falsely cross the threshold, ensuring unconditional exactness for all instances.


\textbf{Theorem 1.} \textit{Algorithm~\ref{alg:hybrid} returns the optimal solution to the 0/1 knapsack problem with probability at least $1 - 2^{-r}$ after $r$ independent repetitions of the quantum detection phase, for all instances.}

\textit{Proof.} We show that no optimal item is permanently excluded by the filtering phase.

Consider any optimal solution $x^*$ and any item $j$ with $x_j^* = 1$. There are three cases:
\begin{enumerate}
    \item If $j \in F_1$, it is included in the solution during Phase~1.
    \item If $j \in C$, it is considered in Phase~2, which searches the B\&B tree over all $2^m$ subsets of the core.
    \item The case $j \in F_0$ cannot occur: by the definition of reduced-cost fixing (Algorithm~\ref{alg:filter}), if $j \in F_0$ then forcing $j$ into the solution causes the LP upper bound to drop below the greedy lower bound $Z_{\text{LB}} \leq Z^*$. Since $Z^*$ is the optimal value and the LP bound is an upper bound, this implies no optimal solution includes $j$.
\end{enumerate}

For case (2), the algorithm tests thresholds via a classical binary-search outer loop. At each threshold, it invokes Montanaro's quantum detection primitive, which correctly identifies if a valid assignment exceeding the threshold exists with high probability. By bounding the failure probability at each of the $O(\log V_{\max})$ steps, the binary search converges to the optimal core assignment with high probability.

Since $F_1$ items are provably optimal, $F_0$ items are provably non-optimal, and the quantum phase finds the optimal core assignment, the combined solution $F_1 \cup x^*_C$ is globally optimal. Repeated applications with independent randomness yield success probability $1 - 2^{-r}$ after $r$ independent runs. \hfill$\blacksquare$

\textbf{Lemma 1b (Monotonicity of the Marginal-Loss Formula).} \label{lem:monotonicity} \textit{For any item $i \ne s$, the true perturbed LP bound satisfies $Z'_i \le \hat{Z}'_i$, where $\hat{Z}'_i$ is the derived formula assuming no split shift ($\hat{Z}'_i = Z_{LP} - w_i |e_i - e_s|$). The formula provides an upper bound on the true perturbed value, and equality holds when the no-split-shift condition is satisfied.}

\textbf{Theorem 2 (Core Containment).} \textit{Let $C_{RC}$ be the reduced-cost core produced by Algorithm~\ref{alg:hybrid}, and let $C_{\delta^*}$ be the $\delta$-band core from Definition~3 with $\delta^* = \frac{Z_{LP} - Z_{LB}}{w_{\min} \cdot e_s}$, where $w_{\min} = \min_{i} w_i$. Assume items are sorted in strictly decreasing efficiency order. Then $C_{RC} \subseteq C_{\delta^*}$.}

\textit{Proof.} Let $i \in C_{RC}$, meaning item $i$ was not fixed and $Z'_i \ge Z_{LB}$. By Lemma 1b, $Z'_i \le \hat{Z}'_i = Z_{LP} - w_i |e_i - e_s|$. Thus, $Z_{LB} \le Z_{LP} - w_i |e_i - e_s|$, which rearranges to $w_i |e_i - e_s| \le Z_{LP} - Z_{LB}$, and therefore $|e_i - e_s| \le \frac{Z_{LP} - Z_{LB}}{w_i} \le \delta^* \cdot e_s$. This implies $i \in C_{\delta^*}$. $\square$

\textbf{Remark.} The no-split-shift condition is not required as a hypothesis for Theorem 2. While violating it breaks the formula's equivalence, Lemma 1b shows it only makes the true $Z'_i$ smaller, which can only cause more items to be fixed, preserving the containment direction. The proof also establishes a tighter item-dependent bound $|e_i - e_s| \le (Z_{LP} - Z_{LB})/w_i$.

\textbf{Proposition 2 (Composed Counting Bound).} \textit{Assume the hypotheses of Theorem 2, and that there exists a minimum efficiency spacing $\Delta > 0$ such that $|e_i - e_j| \ge \Delta$ for all $i \ne j$. Then the core size is bounded by $|C_{RC}| \le \lfloor \frac{2(Z_{LP} - Z_{LB})}{w_{\min} \cdot \Delta} \rfloor + 1$.}

\textit{Proof.} By Theorem 2, $|C_{RC}| \le |C_{\delta^*}|$. The band $C_{\delta^*}$ has width $2 \delta^* e_s$. With spacing $\Delta$, this interval contains at most $\lfloor 2\delta^* e_s / \Delta \rfloor + 1$ items. $\square$

\textbf{Consequence for $T_{LP}$ and Verification.} The LP-pruned B\&B tree on the core satisfies $T_{LP} \le 2^{|C_{RC}|}$. Combined with Montanaro's quantum B\&B, this yields a quantum query complexity bounded by $O(2^{(Z_{LP} - Z_{LB})/(w_{\min} \Delta)} \cdot m \log m \log V_{\max})$. However, there are three important caveats: (1) The $\Delta$-spacing is a structural hypothesis not guaranteed universally (e.g., strongly correlated instances cluster efficiencies, shrinking $\Delta$); (2) $w_{\min}$ can be small; and (3) $2^m$ is a worst-case tree bound, whereas the actual tree is typically smaller.
The containment theorem and Monotonicity Lemma were tested via a rational arithmetic test over 2,000 instances (16,518 item-level checks), and independently reproduced across 3,000 instances (20,339 checks) with zero containment violations observed.

\textbf{Remark (Necessity of $\Delta$).} The $\Delta$-spacing hypothesis in Proposition~2 is \emph{necessary}, not merely a convenient simplification. Consider the family $\mathcal{F}_n$ with $n$ items, all weights $w_i = 1$, values $v_i = 2 - i/n^2$ ($i = 1, \ldots, n$), and capacity $W = \lfloor n/2 \rfloor + \tfrac{1}{2}$. The LP gap converges to $\tfrac{1}{2}v_s \to 1$ (bounded, independent of $n$), while $|C_{RC}| = n$ for every $n \ge 2$---the full item set is in the core and no item is fixed by Algorithm~1. This is not an artefact of the marginal-loss formula: in $\mathcal{F}_n$ every item violates the no-split-shift condition ($w_i = 1 > w_s - R = \tfrac{1}{2}$), so the analysis proceeds from the exact LP re-optimisation directly. When item $i < s$ is forced out, the freed unit of capacity first completes item $s$ (at cost $\tfrac{1}{2}e_s$) then fills item $s{+}1$ fractionally (at cost $\tfrac{1}{2}e_{s+1}$); the true marginal loss is $v_i - \tfrac{1}{2}(v_s + v_{s+1}) = O(1/n^2)$, far below the $O(1)$ gap, so the item survives. The mechanism is that all efficiencies are packed into a window of width $O(1/n^2)$ around $e_s$; split-shifts always land on a neighbour with nearly identical efficiency, keeping every marginal loss negligible relative to the gap. The minimum spacing in $\mathcal{F}_n$ is $\Delta_n = 1/n^2 \to 0$, and Proposition~2's bound grows as $O(n^2)$---correctly diverging, consistent with (but looser than) the true $|C_{RC}| = n$. Hence no bound of the form $|C_{RC}| \le f(Z_{LP} - Z_{LB})$ depending only on the LP gap can exist; the $\Delta > 0$ hypothesis cannot be removed in general.

\textbf{Proposition 1 (Composed Query Complexity).} \textit{The implemented threshold-search composition algorithm solves the 0/1 knapsack problem using a classical outer loop to test thresholds. This implementation composes a classical binary-search outer loop with the fixed-threshold detection primitive, requiring at most $O(\log V_{\max})$ fixed-threshold detection calls. If each detection call evaluates a tree of size $T_{\tau_i}$, bounded by Montanaro's primitive taking $O(\sqrt{T_{\tau_i}}\,m\log m)$ queries, the total implemented algorithm query bound is}
\begin{equation}
    O\!\left(\sqrt{T_{\max}}\,m\log m\log V_{\max}\right) \text{ queries},
    \label{eq:totalcomplexity}
\end{equation}
\textit{where $T_{\max} = \max_i T_{\tau_i}$ over the thresholds tested by the classical binary-search outer loop.}

\textit{Proof.} Phase~1 sorts the $n$ items by efficiency in $O(n \log n)$ time. Computing the greedy lower bound and the initial LP upper bound takes $O(n)$. Phase~2 invokes the threshold detection primitive up to $O(\log V_{\max})$ times. Montanaro's fixed-threshold detection result is the primitive used at each threshold; Theorem 4 is not directly implemented by quantum\_bb\_optimize. Summing at most $O(\log V_{\max})$ fixed-threshold detection calls, each bounded by $O(\sqrt{T_{\tau_i}}\,m\log m)$, gives the stated query complexity. \hfill$\blacksquare$

\subsection{Discussion: Structural Interpretation}
The experiments suggest that the hybrid algorithm adapts across instance distributions. For uncorrelated uniform instances, classical B\&B evaluates relatively few states. In the tested instances, subset-sum capacities generate much larger search spaces. The theorem-backed query bound indicates that, conditional on a large LP-pruned tree, square-root dependence on tree size can reduce query complexity relative to classical traversal of the same tree.

\subsection{Comparison with Prior Work}

Table~\ref{tab:comparison} summarizes the key differences.

\begin{table}[htbp] \caption{Comparison of Knapsack/Subset-Sum Algorithms} \label{tab:comparison} \begin{center} \begin{tabular}{lcccc} \toprule \textbf{Algorithm} & \textbf{Reported Complexity Measure} & \textbf{Space} & \textbf{Prov.} & \textbf{QRAQM} \\ \midrule Horowitz--Sahni & $O(2^{n/2})$ & $O(2^{n/2})$ & Yes & N/A \\ Grover & $O(2^{n/2})$ & $O(n)$ & Yes & No \\ Bernstein et al. & $O(2^{0.241n})$ & $O(2^{0.241n})$ & No$^\dagger$ & Yes \\ Helm \& May & $O(2^{0.226n})$ & $O(2^{0.226n})$ & No$^\dagger$ & Yes \\ Montanaro B\&B & $\tilde{O}(\sqrt{T}\,d)$ queries & $O(\text{poly})$ & Yes & No \\ QTG+AA & $O(\sqrt{N_f})$ & $O(\text{poly})$ & Yes & No \\ \textbf{Ours} & $O(\sqrt{T_{\max}}\,m\log m\log V_{\max})$ & $O(\text{poly})$ & \textbf{Yes} & \textbf{No} \\ \bottomrule \multicolumn{5}{l}{\footnotesize $^\dagger$Relies on Heuristic 2, shown unsound by \cite{bonnetain2020improved}.} \\ \multicolumn{5}{l}{\footnotesize $T$: full B\&B tree. $T_{LP}$: LP-pruned tree ($T_{LP} \ll T$). $N_f$: feasible states.} \end{tabular} \end{center} \end{table}

Our algorithm achieves $O\!\left(\sqrt{T_{\max}}\,m\log m\log V_{\max}\right)$ queries, combining the structural advantage of exact LP reduced-cost fixing with composed fixed-threshold quantum detection.




\section{Experimental Results} \label{sec:experiments}

\subsection{Experimental Setup}
We implemented and validated all three phases of Algorithm~\ref{alg:hybrid} in Python. Crucially, unlike a complexity-counting argument, the explicit Belovs--Montanaro walk operator is constructed computationally for completed trees and spectral/detection conditions are evaluated. The full optimization algorithm \texttt{quantum\_bb\_optimize} uses a classical binary-search outer loop over thresholds $\tau_i$, invoking the Montanaro detection primitive at each threshold. Therefore, the implementation does not directly implement Montanaro Theorem 4. Phases~1 and~3 are exact classical preprocessing/reconstruction. All reported optima are verified against an exact dynamic-programming solve.

The spectral audit tests: the effective spectral-gap inequality, marked phase-zero weight, and the analytical phase-estimation acceptance criterion. It does NOT empirically measure the asymptotic optimization query complexity.

Experiments used five standard instance classes from Pisinger's classification~\cite{ref3}: (1)~uncorrelated, (2)~weakly correlated, (3)~strongly correlated, (4)~subset-sum, and (5)~inverse strongly correlated. For each class, we tested $n \in \{12, 14, 16, 18, 20, 22, 24\}$ with 20 random seeds per configuration.

\subsection{Variable Fixing Results}
Table~\ref{tab:fixing} shows the filtering ratio $\alpha$ across instance types.

\begin{table}[htbp] \caption{Average Core Ratio $\alpha = m/n$ by Instance Type} \label{tab:fixing} \begin{center} \begin{tabular}{lcccc} \toprule \textbf{Instance Type} & $\boldsymbol{\alpha}$ & $|F_1|/n$ & $|F_0|/n$ & \textbf{LP Gap} \\ \midrule Uncorrelated & 0.32 & 0.35 & 0.33 & Tight \\ Weakly correlated & 0.76 & 0.12 & 0.12 & Medium \\ Strongly correlated & 0.95 & 0.02 & 0.03 & Loose \\ Subset-sum & 1.00 & 0.00 & 0.00 & Very loose \\ Inverse strongly & 0.81 & 0.10 & 0.09 & Medium \\ \bottomrule \end{tabular} \end{center} \end{table}

\subsection{Spectral Validation of the Quantum Walk}
Across the 1200 theoretically attempted conditions, 889 successfully completed within the simulable regime ($T_{\text{LP}} \le 2000$). We restrict our spectral conclusions strictly to these 889 simulable trees. The taxonomy is: 489 completed marked, 400 completed unmarked, 74 tree-cap marked, 74 seed-level empty-core conditions (37 bypassed, 37 empty), and 163 immediately pruned unmarked conditions. The selection-bias audit confirms that the pre-eigendecomposition core size $m$ is significantly larger for capped trees (median $m=16$) than for completed trees (median $m=10$).

For the completed instances, we provide computational evidence for three primary spectral endpoints, bounded by exact Clopper-Pearson 95\% confidence intervals (CI):
\begin{enumerate}
    \item \textbf{Unmarked Effective Spectral-Gap Inequality:} satisfied for 400/400 completed unmarked trees (95\% CI: [0.9908, 1.0000]).
    \item \textbf{Marked Phase-Zero Weight:} satisfied for 489/489 completed marked trees (95\% CI: [0.9925, 1.0000]).
    \item \textbf{Phase-Estimation Separation:} 889/889 analytical detection criterion agreement with known marked/unmarked status (95\% CI: [0.9959, 1.0000]).
\end{enumerate}

The analytical phase-estimation acceptance probabilities exhibit strong distributional separation between marked and unmarked trees. We performed an independent-sample Mann-Whitney U test on the 889 completed simulated trees ($T_{\text{LP}} \le 2000$). The marked acceptance probabilities (N=489, median=0.5177, IQR=0.1752) were significantly higher than the unmarked probabilities (N=400, median=0.0311, IQR=0.0999), yielding $U = 195600.0$ and $p = 1.05\times 10^{-145}$.

This Mann-Whitney test supports distributional separation of the analytical phase-estimation acceptance probabilities in the completed simulated subset. It does not support quantum speedup, runtime improvement, quantum advantage, or complexity improvement. The Clopper-Pearson intervals describe the observed satisfaction/agreement proportions in the completed simulable subset, and do not prove universal theorem validity.

\begin{table}[htbp] \caption{Taxonomy of Spectral Validation} \label{tab:measured} \begin{center} \begin{tabular}{lc} \toprule \textbf{Category} & \textbf{Count} \\ \midrule Completed Marked (Phase Zero Validated) & 489 \\ Completed Unmarked (Spectral Gap Validated) & 400 \\ \midrule \textbf{Subtotal Simulable Trees} & \textbf{889 / 889} \\ \midrule Tree-Cap Marked (Truncated $T_{\text{LP}} > 2000$) & 74 \\ Unmarked (Immediately Pruned) & 163 \\ Seed-Level Empty Core & 74 \\ \midrule \textbf{Total Attempted Conditions} & \textbf{1200} \\ \bottomrule \multicolumn{2}{p{\columnwidth}}{\footnotesize Validated against theoretical detection criterion.} \end{tabular} \end{center} \end{table}

\subsection{Classical Runtime and Memory Characterization}
We benchmarked the wall-clock execution time and peak memory consumption of classical B\&B against LP-pruned classical B\&B on a standard Python 3 execution environment. Instrumentation used \texttt{time.perf\_counter\_ns()} and \texttt{tracemalloc}. Across 1500 benchmark records at scales $n \in \{12, 14, 16, 18, 20\}$, we performed a paired Wilcoxon signed-rank test on identical instances. 

For these extremely small scales, classical B\&B without exact LP preprocessing was universally faster in wall-clock time across all classes (paired median runtime ratio $C/LP \approx 0.24$ to $0.62$, Wilcoxon $p < 10^{-4}$ for most). This reflects the constant-factor overhead of repeated Dantzig bound evaluations in Python dominating the trivial search space at $n \le 20$. LP-pruned peak memory remained minuscule (median $1.2\text{--}1.8$ KB). 

The hybrid state-vector simulation runtime is not an estimate of physical quantum execution time and is therefore excluded from claims of wall-clock quantum acceleration. We do not report any wall-clock quantum speedup.

\subsection{Comparison Against Google OR-Tools at Scale}

We benchmarked the algorithm against Google OR-Tools (version 9.15.6755) across 5 instance classes, with 5 trials each, scaling $n$ up to 1000 for tractable classes and hard-stopping at $n=500$ for strongly and inverse-strongly correlated classes once intractability was reached. Across every class and size, Phase 1 (reduced-cost fixing) resolved 0 of the tested instances outright; Phase 2 core-solving was required in 100\% of cases tested, providing large-scale context for the variable fixing effectiveness reported for smaller $n$ in Section VI-B (Table~\ref{tab:fixing}). Of 90 total instances attempted across all classes and sizes, 69 were fully solved to certified optimality by both solvers, with zero mismatches between this paper's reported optimal value and OR-Tools' optimal value in every one of those 69 cases; this supports the correctness of our implementation. 

For the remaining 21 instances, cutoffs were reached (OR-Tools at 20s, our B\&B at $1{,}000{,}000$ nodes); because cutoff time is not solve time, no timing comparison is drawn from these rows, and Table~\ref{tab:ortools_cat_b} reports only the timeout counts. Table~\ref{tab:ortools_cat_a} details the performance for the 69 fully solved instances.

\begin{table}[htbp]
\caption{Category (a): Both Solvers Reached Certified Optimal}
\label{tab:ortools_cat_a}
\begin{center}
\begin{tabular}{lcccccc}
\toprule
\textbf{Instance Class} & $\boldsymbol{n}$ & \textbf{Successes} & \textbf{Core $m$} & \textbf{OR-Tools (s)} & \textbf{Our B\&B (s)} & \textbf{Notes} \\
\midrule
Uncorrelated & 100 & 5/5 & 18.2 & 0.0001 & 0.0031 & \\
Uncorrelated & 200 & 5/5 & 27.2 & 0.0000 & 0.0210 & \\
Uncorrelated & 500 & 5/5 & 78.0 & 0.0002 & 0.0984 & \\
Uncorrelated & 1000 & 5/5 & 75.6 & 0.0001 & 0.3301 & \\
Weakly correlated & 100 & 5/5 & 79.4 & 0.0000 & 0.0032 & \\
Weakly correlated & 200 & 5/5 & 56.4 & 0.0002 & 0.0107 & \\
Weakly correlated & 500 & 5/5 & 91.6 & 0.0001 & 0.0729 & \\
Weakly correlated & 1000 & 5/5 & 126.2 & 0.0001 & 0.3253 & \\
Subset sum & 100 & 5/5 & 100.0 & 0.0000 & 0.0028 & \\
Subset sum & 200 & 5/5 & 200.0 & 0.0000 & 0.0109 & \\
Subset sum & 500 & 5/5 & 500.0 & 0.0000 & 0.0733 & \\
Subset sum & 1000 & 5/5 & 1000.0 & 0.0001 & 0.2970 & \\
Strongly correlated & 100 & 3/5 & 100.0 & 0.0809 & 0.0877 & \\
Strongly correlated & 200 & 2/5 & 200.0 & 0.0130 & 0.0607 & Warning: $n=2$, not a stable avg \\
Inverse strongly & 100 & 3/5 & 36.0 & 0.0563 & 0.0191 & \\
Inverse strongly & 200 & 1/5 & 25.0 & 0.0017 & 0.0108 & Warning: $n=1$, not a stable avg \\
\bottomrule
\end{tabular}
\end{center}
\end{table}

\begin{table}[htbp]
\caption{Category (b): Cutoffs Reached}
\label{tab:ortools_cat_b}
\begin{center}
\begin{tabular}{lccc}
\toprule
\textbf{Instance Class} & $\boldsymbol{n}$ & \textbf{OR-Tools Timeouts} & \textbf{Our B\&B Timeouts} \\
\midrule
Strongly correlated & 100 & 2/5 & 2/5 \\
Strongly correlated & 200 & 3/5 & 3/5 \\
Strongly correlated & 500 & 5/5 & 5/5 \\
Inverse strongly & 100 & 2/5 & 2/5 \\
Inverse strongly & 200 & 4/5 & 4/5 \\
Inverse strongly & 500 & 5/5 & 5/5 \\
\bottomrule
\end{tabular}
\end{center}
\end{table}


\subsection{Empirical Scaling of $T_{\text{LP}}$ with LP Gap and Core Size}

We empirically characterize the scaling of $T_{\text{LP}}$ with instance parameters using a regression model fit to our benchmark dataset. We fit the model $\log(T_{\text{LP}}) \sim \log(\text{gap}) + m$, where gap is the true integrality gap $Z_{\text{LP}} - Z^*$ computed against the exact optimal value, and $m$ is the core size. The pooled model yields $R^2 = 0.469$ (5-fold CV $R^2 = 0.460$), but this masks substantial class dependency.

\begin{table}[htbp]
\caption{Per-Class Regression Fits for $\log(T_{\text{LP}}) \sim \log(\text{gap}) + m$}
\label{tab:tlp_regression}
\begin{center}
\begin{tabular}{lrr}
\toprule
\textbf{Instance Class} & $\boldsymbol{R^2}$ & \textbf{CV $\boldsymbol{R^2}$} \\
\midrule
Uncorrelated & 0.627 & 0.584 \\
Weakly correlated & 0.594 & 0.479 \\
Strongly correlated & 0.601 & 0.560 \\
Subset sum & 0.010 & -0.074 \\
Inverse strongly & 0.583 & 0.540 \\
\bottomrule
\end{tabular}
\end{center}
\end{table}

The coefficient for $\log(\text{gap})$ is $\approx 1.01$ and highly significant ($p < 0.001$) for strongly correlated instances, but statistically indistinguishable from zero ($p = 0.919$) for uncorrelated instances, and mathematically undefined for subset sum (where the gap is identically zero for all instances). This behavior is consistent with the established structural properties of these instances (as seen in Table~\ref{tab:fixing}). Core size $m$, however, is a consistent positive predictor across all valid classes (coefficients $\approx 0.11\text{--}0.13$). We do not claim any universal closed-form bound on $T_{\text{LP}}$; this is a fitted empirical trend specific to the tested instance distributions and parameter ranges ($n \le 24$).

\section{Discussion}

\subsection{Instance-Dependent Query-Complexity Behavior}

Our results reveal a key structural insight: the sources of search reduction are instance-dependent and complementary, creating a synergistic algorithm that exploits the strengths of both classical preprocessing and quantum search.

For \textit{uncorrelated instances}, the LP relaxation bound is tight (LP gap $\approx O(1)$), enabling aggressive variable fixing ($\alpha \approx 0.32$ at $n = 24$, dropping to median $\alpha = 0.0276$ at $n = 10{,}000$). The B\&B tree is correspondingly small ($T_{\text{LP}} \approx 70$ at $n=24$). Here, most of the benefit comes from Phase~1 (classical), with the quantum phase contributing a modest additional $\sqrt{T_{\max}}\,m\log m / T_{\max}$ factor in query count.

For \textit{strongly correlated and subset-sum instances}, variable fixing is ineffective ($\alpha \to 1$) because all item efficiencies are nearly identical. However, the tested strongly correlated and subset-sum instances produced substantially larger LP-pruned trees ($T_{\text{LP}} \approx 3079$ at $n=24$). Here, Phase~2 provides the key query reduction: Montanaro's quantum walk achieves a quantum query reduction over substantially larger LP-pruned B\&B trees.

This synergy is the core novelty: \textit{In the tested instance classes, the relative contribution of LP fixing and the theorem-backed quantum query bound varies substantially with instance structure, and the two phases address different algorithmic bottlenecks}. For uncorrelated instances, $\alpha \ll 1$ makes the quantum phase trivial; For larger LP-pruned trees, the fixed-threshold detection bound has square-root dependence on tree size, with the explicit depth factor $m$. In both regimes, the algorithm adapts automatically.

\subsection{Hardware Validation Insights}

Our IBM \texttt{ibm\_fez} experiments reveal a clear NISQ noise boundary for amplitude amplification circuits. Table~\ref{tab:hardware} summarizes the measured hardware success probability versus the ideal simulator success probability across scaled core sizes.

\begin{table}[htbp]
\caption{IBM Hardware Results vs. Simulator Success Probability}
\label{tab:hardware}
\begin{center}
\footnotesize
\begin{tabular}{lcccc}
\toprule
\textbf{Instance} & \textbf{Sim Prob} & \textbf{HW Prob} & \textbf{Transpiled Depth} & \textbf{2Q Gates} \\
\midrule
Core-3A (Easy) & 0.9453 & 0.6743 & 301 & 78 \\
Core-3B (Medium) & 0.9453 & 0.6671 & 294 & 72 \\
Core-4A (Hard) & 0.9613 & 0.1235 & 2106 & 661 \\
Core-4B (Subset-Sum) & 0.9613 & 0.1344 & 2102 & 659 \\
Core-5A (Uncorrelated) & 0.9992 & 0.0237 & 12512 & 4017 \\
Core-5B (Correlated) & 0.9613 & 0.0637 & 9273 & 3012 \\
Core-6A (Hard) & 0.9992 & 0.0312 & 50520 & 16579 \\
Core-6B (Subset-Sum) & 0.9981 & 0.0330 & 38129 & 12456 \\
\bottomrule
\end{tabular}
\end{center}
\end{table}

At $n \leq 4$ (circuit depth $\leq 2{,}100$, $\leq 661$ two-qubit gates), the correct optimal solution emerges as the dominant measurement outcome with success probabilities of 12--67\%. At $n \geq 5$ (circuit depth $\geq 9{,}000$, $\geq 3{,}000$ two-qubit gates), noise overwhelms the quantum signal, and the output approaches a uniform distribution.

Fidelity was not independently measured. Backend calibration snapshots were not retained as controlled experimental covariates; consequently, execution-time T1, T2, readout-error, and two-qubit-error statistics are not used for quantitative attribution of the observed hardware--simulator discrepancy.

This characterization has practical implications. For our hybrid algorithm, the quantum phase operates on the \textit{core} items only ($m = \alpha n$). Since $\alpha \ll 1$ for uncorrelated and weakly correlated instances, LP fixing reduces the number of variables exposed to the quantum phase, but the present hardware experiments support only very small cores. No claim is made that current hardware can execute the full quantum B\&B phase for large reduced cores.


\subsection{Limitations}

\textbf{Hardware requirements.} The fixed-threshold detection primitive has an $O(\sqrt{T_{\tau}}\,m\log m)$ query bound under the Montanaro oracle-model assumptions, each involving an LP bounding oracle with reversible marked and threshold predicates. For large core sizes ($m \geq 5$), this demands fault-tolerant quantum hardware, as our IBM experiments demonstrate.

\textbf{NISQ noise boundary.} Our hardware experiments show that current Heron~r2 processors can reliably execute amplitude amplification only for cores with $n \leq 4$ using our unitary oracle construction. Decomposed gate-level oracles and error mitigation techniques~\cite{ref31} may extend this boundary.

\textbf{Small-$n$ overhead.} At $n \leq 24$, the Montanaro $\text{poly}(d)$ overhead can make the quantum algorithm slower than classical B\&B for easy instances. The quantum query reduction is asymptotic and becomes pronounced as $n$ grows and the LP-pruned tree becomes sufficiently large.

\subsubsection{Resource-Estimation Crossover Analysis}
This analysis compares the quantum T-gate cost of the composed query bound against the classical per-node operation cost to derive the crossover threshold. This is strictly an operation-count comparison, not a wall-clock time estimate. Table~\ref{tab:crossover} details the combined per-query T-cost ($C_Q$), aggregating the resources for predicate $V$ (Table~\ref{tab:oracle}) and the LP heuristic $h$ (Table~\ref{tab:h_resources}). For $m=4$, this includes the newly measured T-count for $h$, while for $m=6, 8$, we explicitly note that these counts remain a partial lower bound because full gate-level decomposition is deferred~\cite{vedral1996}.

\begin{table}[htbp]
\caption{Per-Query T-Cost Estimates ($C_Q$) and Required Crossover Thresholds ($T_{\max}$)}
\label{tab:crossover}
\begin{center}
\begin{tabular}{lrr}
\toprule
\textbf{Core Size ($m$)} & \textbf{$C_Q$} & \textbf{Required $T_{\max}$ Crossover} \\
\midrule
4$^\S$ & $\approx 1{,}388{,}704$ & $\approx 1.88 \times 10^{15}$ \\
6 & $\ge 240{,}986$ & $\ge 3.49 \times 10^{14}$ \\
8 & $\ge 269{,}760$ & $\ge 1.50 \times 10^{15}$ \\
\bottomrule
\multicolumn{3}{l}{\footnotesize $^\S$MEASURED: $\approx$ denotes a synthesized gate count (Table~\ref{tab:h_resources}). $\ge$ rows remain partial lower bounds (gate-level decomposition deferred).}
\end{tabular}
\end{center}
\end{table}

Setting the classical cost ($T_{\max}$) equal to the quantum cost gives the crossover condition: $T_{\max} \ge (C_Q \cdot m \log_2 m \cdot \log_2 V_{\max})^2$. An adversarial comparison against the empirical data reveals that none of the tested instances approach this required crossover threshold by many orders of magnitude; the maximum observed tree sizes ($T_{\text{LP}} \approx 70$ for uncorrelated and $T_{\text{LP}} \approx 3{,}079$ for strongly correlated instances at $n=24$) are ten to eleven orders of magnitude too small. Because the T-cost estimate is a strict lower bound, the true crossover point can only be larger, strengthening the conclusion that the overhead of reversible arithmetic overwhelms the quantum query reduction for the evaluated regime.

\textbf{Pseudo-polynomial regime.} For instances with small integer coefficients, dynamic programming runs in $O(nW)$ time and space, and dominates: our theoretical comparison to branch-and-bound is meaningful only in the large-coefficient regime where DP is infeasible due to time and memory limits. We make this scope explicit rather than claiming general advantage over all classical methods.

\textbf{Detection reliability and simulability.} Our phase-estimation precision constant is calibrated empirically. Walk simulation restricts the spectral audit to trees with $T_{\text{LP}} \le 2000$ nodes. \textbf{Oracle Completeness.} Proposition~1 implements the reversible evaluation/threshold predicate $V$. For the LP bounding heuristic $h$, we have specified and verified an exact fixed-precision arithmetic formula (Lemma \ref{lem:precision}) and established the complete register layout and qubit resource estimates (Table~\ref{tab:h_resources}). However, full gate-level decomposition of the required arithmetic blocks (multipliers and dividers) remains deferred. Consequently, while the mathematical exactness and qubit overhead of $h$ are confirmed, we still do not claim a full-algorithm gate-complexity bound.




\section{Conclusion}

We presented a hybrid algorithm that couples LP-guided reduced-cost variable fixing with Montanaro's quantum branch-and-bound for the 0/1 knapsack problem---to our knowledge among the first to combine provability, gate-model feasibility, knapsack-specific LP fixing, instance adaptivity, and computational spectral validation of the explicit walk operator (Table~\ref{tab:related}). The implemented threshold-search composition has query complexity $O\!\left(\sqrt{T_{\max}}\,m\log m\log V_{\max}\right)$, where $T_{\max}$ is the maximum LP-pruned tree size encountered and $m = \alpha n$ is the core size after variable fixing. We established two key results:

\begin{enumerate}     \item \textit{Correctness}: reduced-cost fixing is exact---no heuristic assumptions are involved---and Montanaro's quantum walk finds the optimal B\&B solution with high probability.     \item \textit{Instance-dependent synergy}: for uncorrelated instances, LP fixing achieves median $\alpha = 0.0276$ at scale ($n = 10{,}000$), reducing the effective search to a trivial core; for the tested strongly constrained instances with substantially larger LP-pruned trees, the theorem-backed fixed-tree detection primitive has square-root dependence on tree size, up to the stated depth factors, relative to classical traversal of the same tree.     \end{enumerate}

Experiments on 700 instances across five Pisinger difficulty classes confirmed 100\% correctness against exact dynamic programming. Critically, a computational simulation of the quantum walk tested the underlying effective spectral-gap inequality and phase-zero weight constraints across 889 completed simulated trees ($T_{\text{LP}} \le 2000$), providing computational evidence that the implemented LP-pruned walk satisfies the tested spectral conditions underlying Montanaro's detection analysis. A fault-tolerant resource estimate for our answer-agnostic reversible oracle establishes baseline component costs for the predicate $V$, and a NISQ characterization on IBM's \texttt{ibm\_fez} processor (156-qubit Heron~r2) identified the correct solution as the top-1 measured state for $n \leq 4$ while delineating the noise boundary at $n \geq 5$.

Future work includes: (a)~extending the analysis to multi-dimensional and multi-objective knapsack variants; (b)~implementing decomposed gate-level oracles to reduce circuit depth and extend the hardware-viable range beyond $n = 4$ core items; (c)~integrating quantum error mitigation techniques to improve hardware success probabilities at larger core sizes; (d)~characterizing instance classes or realistic distributions under which the minimum efficiency spacing $\Delta$ remains bounded away from zero---a $\Delta$-free analytical bound on $|C_{RC}|$ in terms of the LP gap alone is ruled out by the necessity construction in Remark~(Necessity of $\Delta$), so the open question is to identify structured families admitting useful $\Delta$ lower bounds or to derive tighter bounds for specific distributions such as uncorrelated or weakly correlated instances; and (e)~evaluating the algorithm on next-generation fault-tolerant quantum processors as they become available.




\section{Data Availability and Reproducibility}

To ensure full reproducibility and scientific transparency, the complete source code for our classical B\&B, quantum walk simulation, oracle compilation, and IBM Quantum hardware experiments is publicly available. The repository includes scripts to regenerate the 700 Pisinger instances, verify all optimizations against dynamic programming, and reproduce all tables and figures. Crucially, the repository simulates the exact Belovs--Montanaro walk operator computationally to test spectral conditions, moving beyond theoretical complexity assertions.










\bibliographystyle{IEEEtran} \bibliography{references}

\end{document} 